% GENERATED FILE. Edit canonical lessons, then run npm run generate:books.
% canonical-source-hash: fnv1a64:52aad6b65f9a0c31
% canonical-lessons: MW-C226-tibo, MW-C226-ucho, MW-C226-suno, MW-C226-thal, MW-C226-ujar

\chapter{Sand Hills}
\label{ch:mw-sand-hills}
\hlchaptermodality{226}

\begin{chapteropening}
\textbf{By the end of this chapter:} \emph{I can describe the sand hills and the empty land.}

Complete the last lesson of chapter 226: I can describe the sand hills and the empty land.
\end{chapteropening}

\section[\texorpdfstring{\mw{टीबो} (ṭībo)}{टीबो (ṭībo)}]{\mw{टीबो} (ṭībo) — a sand hill; a dune}

\label{lesson:MW-C226-tibo}



\begin{quote}
Before the new one: say the Marwadi for bedding, then the Marwadi for a woollen blanket.
\end{quote}

\subsection*{You'll want to know: \mw{टीबो}}
\textbf{\mw{टीबो}} — \emph{ṭībo} — \textquotedblleft{}a sand hill; a dune\textquotedblright{}.

\begin{quote}
\mw{गांव} \mw{कनै} \mw{एक} \mw{टीबो} \mw{है।}
\end{quote}

\begin{quote}
\emph{gā̃v kanai ek ṭībo hai.}
\end{quote}

\begin{quote}
There's a sand hill near the village.
\end{quote}

A masculine noun.

\begin{cousinweb}
\textbf{\mw{टीबो}} (\emph{ṭībo}) is inherited from a reconstructed \emph{*ṭappa}, \textquotedblleft{}a hill\textquotedblright{}, one of several shapes CDIAL gives for it; CDIAL lists Hindi \emph{ṭībā}, \textquotedblleft{}a hill, a sandbank\textquotedblright{}, and a Gujarati cousin. CDIAL gives no Marwari form in this entry (CDIAL 5446).
\end{cousinweb}

\begin{grammarlens}[title={What you've built}]
One more word for this chapter.
\end{grammarlens}

\subsection*{Your turn}
\begin{itemize}
\raggedright
  \item \emph{Say it:} \emph{ṭībo}
  \item \emph{Say it:} \emph{ṭībo}, once more
  \item \emph{Say it:} \emph{gā̃v kanai ek ṭībo hai.} — There's a sand hill near the village.
  \item \emph{Recall:} say the Marwadi for bedding, then the Marwadi for a woollen blanket, then say \emph{ṭībo} again
\end{itemize}

\begin{tcolorbox}[breakable,colback=teal!4,colframe=teal!35!black,title={Before you move on}]
What does \mw{टीबो} mean? (\textquotedblleft{}A sand hill; a dune\textquotedblright{}.) Say it once more.
\end{tcolorbox}
\section[\texorpdfstring{\mw{ऊंचो} (ū̃cho)}{ऊंचो (ū̃cho)}]{\mw{ऊंचो} (ū̃cho) — high; tall}

\label{lesson:MW-C226-ucho}



\begin{quote}
Before the new one: say the Marwadi for a woollen blanket, then the Marwadi for a sand hill.
\end{quote}

\subsection*{You'll want to know: \mw{ऊंचो}}
\textbf{\mw{ऊंचो}} — \emph{ū̃cho} — \textquotedblleft{}high; tall\textquotedblright{}.

\begin{quote}
\mw{टीबो} \mw{घणो} \mw{ऊंचो} \mw{है।}
\end{quote}

\begin{quote}
\emph{ṭībo ghaṇo ū̃cho hai.}
\end{quote}

\begin{quote}
The sand hill is very high.
\end{quote}

\begin{cousinweb}
\textbf{\mw{ऊंचो}} (\emph{ū̃cho}) is inherited from Sanskrit \emph{ucca}, \textquotedblleft{}tall, lofty\textquotedblright{}; CDIAL lists Old Marwari \emph{uco}, \textquotedblleft{}high\textquotedblright{} (CDIAL 1634). \textbf{\mw{नीचो}}, low, is its opposite.
\end{cousinweb}

\begin{grammarlens}[title={What you've built}]
One more word for this chapter.
\end{grammarlens}

\subsection*{Your turn}
\begin{itemize}
\raggedright
  \item \emph{Say it:} \emph{ū̃cho}
  \item \emph{Say it:} \emph{ū̃cho}, once more
  \item \emph{Say it:} \emph{ṭībo ghaṇo ū̃cho hai.} — The sand hill is very high.
  \item \emph{Recall:} say the Marwadi for a woollen blanket, then the Marwadi for a sand hill, then say \emph{ū̃cho} again
\end{itemize}

\begin{tcolorbox}[breakable,colback=teal!4,colframe=teal!35!black,title={Before you move on}]
What does \mw{ऊंचो} mean? (\textquotedblleft{}High; tall\textquotedblright{}.) Say it once more.
\end{tcolorbox}
\section[\texorpdfstring{\mw{सूनो} (sūno)}{सूनो (sūno)}]{\mw{सूनो} (sūno) — lonely; deserted}

\label{lesson:MW-C226-suno}



\begin{quote}
Before the new one: say the Marwadi for a sand hill, then the Marwadi for high.
\end{quote}

\subsection*{You'll want to know: \mw{सूनो}}
\textbf{\mw{सूनो}} — \emph{sūno} — \textquotedblleft{}lonely; deserted\textquotedblright{}.

\begin{quote}
\mw{रात} \mw{गेलो} \mw{सूनो} \mw{है।}
\end{quote}

\begin{quote}
\emph{rāt gelo sūno hai.}
\end{quote}

\begin{quote}
At night the path is lonely.
\end{quote}

\begin{cousinweb}
\textbf{\mw{सूनो}} (\emph{sūno}) is inherited from Sanskrit \emph{śūnya}, \textquotedblleft{}empty\textquotedblright{} (CDIAL 12567, which lists Marwari \emph{sūno}). It is said of a place with nobody in it.
\end{cousinweb}

\begin{grammarlens}[title={What you've built}]
One more word for this chapter.
\end{grammarlens}

\subsection*{Your turn}
\begin{itemize}
\raggedright
  \item \emph{Say it:} \emph{sūno}
  \item \emph{Say it:} \emph{sūno}, once more
  \item \emph{Say it:} \emph{rāt gelo sūno hai.} — At night the path is lonely.
  \item \emph{Recall:} say the Marwadi for a sand hill, then the Marwadi for high, then say \emph{sūno} again
\end{itemize}

\begin{tcolorbox}[breakable,colback=teal!4,colframe=teal!35!black,title={Before you move on}]
What does \mw{सूनो} mean? (\textquotedblleft{}Lonely; deserted\textquotedblright{}.) Say it once more.
\end{tcolorbox}
\section[\texorpdfstring{\mw{थळ} (thaḷ)}{थळ (thaḷ)}]{\mw{थळ} (thaḷ) — dry land; the open desert}

\label{lesson:MW-C226-thal}



\begin{quote}
Before the new one: say the Marwadi for high, then the Marwadi for lonely.
\end{quote}

\subsection*{You'll want to know: \mw{थळ}}
\textbf{\mw{थळ}} — \emph{thaḷ} — \textquotedblleft{}dry land; the open desert\textquotedblright{}.

\begin{quote}
\mw{थळ} \mw{मांय} \mw{रेत} \mw{घणी} \mw{है।}
\end{quote}

\begin{quote}
\emph{thaḷ mā̃y ret ghaṇī hai.}
\end{quote}

\begin{quote}
There's a lot of sand in the open desert.
\end{quote}

A masculine noun.

\begin{cousinweb}
\textbf{\mw{थळ}} (\emph{thaḷ}) is inherited from Sanskrit \emph{sthala}, \textquotedblleft{}dry land\textquotedblright{}; CDIAL lists Old Marwari \emph{thala}, \textquotedblleft{}land\textquotedblright{} (CDIAL 13744).
\end{cousinweb}

\begin{grammarlens}[title={What you've built}]
One more word for this chapter.
\end{grammarlens}

\subsection*{Your turn}
\begin{itemize}
\raggedright
  \item \emph{Say it:} \emph{thaḷ}
  \item \emph{Say it:} \emph{thaḷ}, once more
  \item \emph{Say it:} \emph{thaḷ mā̃y ret ghaṇī hai.} — There's a lot of sand in the open desert.
  \item \emph{Recall:} say the Marwadi for high, then the Marwadi for lonely, then say \emph{thaḷ} again
\end{itemize}

\begin{tcolorbox}[breakable,colback=teal!4,colframe=teal!35!black,title={Before you move on}]
What does \mw{थळ} mean? (\textquotedblleft{}Dry land; the open desert\textquotedblright{}.) Say it once more.
\end{tcolorbox}
\section[\texorpdfstring{\mw{उजाड़} (ujāṛ)}{उजाड़ (ujāṛ)}]{\mw{उजाड़} (ujāṛ) — laid waste; desolate}

\label{lesson:MW-C226-ujar}



\begin{quote}
Before the new one: say the Marwadi for lonely, then the Marwadi for dry land.
\end{quote}

\subsection*{You'll want to know: \mw{उजाड़}}
\textbf{\mw{उजाड़}} — \emph{ujāṛ} — \textquotedblleft{}laid waste; desolate\textquotedblright{}.

\begin{quote}
\mw{वो} \mw{गांव} \mw{उजाड़} \mw{है।}
\end{quote}

\begin{quote}
\emph{vo gā̃v ujāṛ hai.}
\end{quote}

\begin{quote}
That village is deserted.
\end{quote}

\begin{cousinweb}
\textbf{\mw{उजाड़}} (\emph{ujāṛ}) is inherited from a reconstructed \emph{*ujjaṭa}, \textquotedblleft{}uprooted, laid waste\textquotedblright{}; CDIAL lists Old Marwari \emph{ujāṛi}, \textquotedblleft{}a deserted place\textquotedblright{} (CDIAL 1660b).
\end{cousinweb}

\begin{grammarlens}[title={What you've built}]
One more word for this chapter.
\end{grammarlens}

\subsection*{Your turn}
\begin{itemize}
\raggedright
  \item \emph{Say it:} \emph{ujāṛ}
  \item \emph{Say it:} \emph{ujāṛ}, once more
  \item \emph{Say it:} \emph{vo gā̃v ujāṛ hai.} — That village is deserted.
  \item \emph{Recall:} say the Marwadi for lonely, then the Marwadi for dry land, then say \emph{ujāṛ} again
\end{itemize}

\begin{tcolorbox}[breakable,colback=teal!4,colframe=teal!35!black,title={Before you move on}]
What does \mw{उजाड़} mean? (\textquotedblleft{}Laid waste; desolate\textquotedblright{}.) Say it once more.
\end{tcolorbox}
